%%%%%%%%%%%%
%
% $Autor: Wings $
% $Datum: 2019-03-05 08:03:15Z $
% $Pfad: Automatisierung/Skript/Produktspezifikation/Powerpoint/AMF.tex $
% $Version: 4250 $
% !TeX spellcheck = en_GB/de_DE
% !TeX encoding = utf8
% !TeX root = filename 
% !TeX TXS-program:bibliography = txs:///biber
%
%%%%%%%%%%%%


\chapter{Gesichtserkennung mit Edge Impulse über die Kommandozeilenschnittstelle CLI}

In diesem Kapitel wird eine Modell zur Gesichtserkennung mit Hilfe der Palttform Edge Impulse\index{Edge Impulse} entwickelt. Das Modell wird als TensorFlow Lite-Modell gespeichert und mit der Hilfe von \ac{cli} auf das Arduino Board gebracht.
Das Flussdiagramm für den Gesamtprozess ist in der folgenden Abbildung \ref{FaceFlowChart} dargestellt. \Mynote{FLussdiagramm überarbeiten oder entfernen.}

\begin{figure}[H]
	\centering
	\includegraphics[width=\textwidth]{Arduino/flowchart2}
	\caption{Flussdiagramm des Gesamtprozesses}
	\label{FaceFlowChart}
\end{figure}


\section{Verbinden des Arduino-Board mit Edge Impulse}

Das Arduino Board kann von Edge Impulse aus angesteuert werden, wenn \ac{cli} genutzt wird. Dazu muss Edge Impulse entsprechend konfiguriert werden.

 Zunächst muss  Python 3 installiert  und Node.js Version 14.17 oder höher von der Website \url{https://nodejs.org/en/} heruntergeladen werden. Die heruntergeladene Installationsdatei für Windows wird als \FILE{node-v14.17.4-x64.msi} bezeichnet. Wenn die Datei geöffnet wird, startet die Installation, sobald die Schaltfläche \glqq Installation\grqq, wie in der Abbildung \ref{InstallationNode.js} gezeigt, betätigt wird. Der Installationspfad könnte \path{C:\ProgrammDaten\Microsoft\Windows\Startmenü\Programme\Node.js} lauten. Nachdem die Installation abgeschlossen ist, muss die Eingabeaufforderung geöffnet werden und
 
 \medskip
 
  \SHELL{npm install -g edge-impulse-cli --force} 
  
  \medskip 
  
eingeben werden, um die Kommandozeilenschnittstelle für Edge Impulse zu installieren.

\begin{figure}[H]
	\centering
	\includegraphics[width=\textwidth]{Arduino/nodejs}
	\caption{Installation von Node.js }	\label{InstallationNode.js}
\end{figure}



Nach der Installation kann das Arduino Board, zum Beispiel Arduino Portenta H7 oder Arduino Nicla Vision, mit Edge Impulse verbunden werden. Um die aktuelle Entwicklung zu nutzen, muss sichergestellt werden, dass die  neueste Edge Impulse-Firmware von der Website von Edge Impulse \cite{EdgeImpulseBoards:2022} heruntergeladen ist. Auf der Website ist dokumentiert, welche Boards aktuell unterstützt werden.
% \url{https://docs.edgeimpulse.com/docs/arduino-portenta-h7} 
Für das Arduino Board Portenta H7 trägt die heruntergeladene Datei den Namen \FILE{arduino-portenta-h7-preview.zip}. Wenn die Datei entpackt wurde, muss das Flash-Skript für Windows geöffnt werden, das sich in dem Ordner \FILE{flash\_windows.bat} befindet. Nach dem Flashen muss  die Taste \glqq Reset\grqq{} einmal gedrückt werden, um die neue Firmware zu starten. Nun muss in einer Eingabeaufforderung das Kommanmdo

\medskip

 \SHELL{edge-impulse-daemon} 
 
 \medskip
 
 eingegeben werden. Daraufhin wird man aufgefordert, sich mit seiner Edge Impulse ID\index{Edge Impulse ID} und seinem Passwort\index{Edge Impulse Password}  in der Eingabeaufforderung anzumelden. Nachdem die Anmeldung erfolgreich war, kann das Gerät benannt werden, hier wird  es mit  \glqq Arduino Portenta H7\grqq{} bezeichnet, wie in der Abbildung \ref{ArduinoBoardEdgeImpulse} gezeigt.

\begin{figure}[H]
	\centering
	\includegraphics[width=\textwidth]{Arduino/edgecli}
	\caption{Verbindung zwischen Arduino Portenta H7 und Edge Impulse }
	\label{ArduinoBoardEdgeImpulse}
\end{figure}


Um zu überprüfen, ob das Board verbunden ist, wird nun die Website von Edge Impulse geöffnet und auf die Schaltfläche \glqq Geräte\grqq{} geklickt. Dort wird das Gerät, wie in der Abbildung \ref{ArduinoConnectedEdgeImpulse} dargestellt, aufgelistet. JEtzt ist das Board vernunden und es kann it dem Training begonnen werden.
\begin{figure}[H]
	\centering
	\includegraphics[width=\textwidth]{Arduino/connectededge}
	\caption{Verbindung eines Arduino-Boards mit Edge Impulse}
	\label{ArduinoConnectedEdgeImpulse}
\end{figure}

%Um mit dem Vision Shield Bilder zu machen, müssen wir unser Gerät mit dem Edge Impulse-Portal verbinden und dann mit der Aufnahme von Bildern beginnen. Dazu müssen wir die folgende Software installieren.




\subsection{Datenbank}  

Als Grundlage des Trainings wird der Datensatz UTKFace, der in Kapitel \ref{DatasetUTKFace} beschrieben wird. Es können eigene Bilder hinzugefügt werden. Da Edge Impulse nur eine beschränkte Trainingszeit pro Tag kostenfrei zur Verfügung stellt, wird eine Auswahl getroffen:

\begin{itemize}
	\item Trainingssatz: Bestehend aus 388 Bildern
	\item Testsatz: bestehend aus 98 Bildern
\end{itemize}


\bigskip

Nachdem der Datensatz gewählt wurde, kann man nun in Edge Impulse das Training dem Kapitel \ref{EdgeImpulse} folgend trainieren.

\section{Prüfung des Modells}

Nach dem Training des Modells kann auch die Live-Klassifizierung mit der Hardware durchgeführt werden. Dazu muss in Edge Impulse im Menü auf die Option \glqq live classification\grqq{} geklickt werden. Dort muss das Arduino Board durch die Anwahl mit Edge Impulse verbunden werden. Sobald das Arduino Board angeschlossen ist, wird der Videostream angezeigt. Wird nun das Kamerabild nun auf  ein Gesicht konzentriert,  in der Abbildung \ref{EdgeIMpulkseLiveClassification} zu sehen ist, so wird eine Vorhersage getroffen. In diesem Beispiel ist die Wahrscheinlichkeit, dass ein Gesicht detektiert wurde, mit einem Wert von 89\% angegeben.
	
	\begin{figure}[H]
		\centering
		\includegraphics[width=\textwidth]{Arduino/liveclassification}
		\caption{Live-Klassifizierung mit Hilfe der Vision-Shield-Kamera}
		\label{EdgeIMpulkseLiveClassification}
	\end{figure}
	
